\documentclass[11pt]{amsart}
\usepackage{amsmath,amssymb,amsthm,mathtools}
\usepackage[margin=1.05in]{geometry}

\numberwithin{equation}{section}
\newtheorem{theorem}{Theorem}[section]
\newtheorem{proposition}[theorem]{Proposition}
\newtheorem{lemma}[theorem]{Lemma}
\newtheorem{corollary}[theorem]{Corollary}
\theoremstyle{definition}
\newtheorem{definition}[theorem]{Definition}
\theoremstyle{remark}
\newtheorem{remark}[theorem]{Remark}
\newtheorem{example}[theorem]{Example}

\newcommand{\Z}{\mathbb{Z}}
\newcommand{\Q}{\mathbb{Q}}
\newcommand{\R}{\mathbb{R}}
\newcommand{\N}{\mathbb{N}}
\newcommand{\C}{\mathbb{C}}
\newcommand{\T}{\mathbb{T}}
\newcommand{\cA}{\mathcal{A}}
\newcommand{\cG}{\mathcal{G}}
\newcommand{\cR}{\mathcal{R}}
\newcommand{\cT}{\mathcal{T}}
\newcommand{\pp}{\mathfrak{p}}
\newcommand{\aaa}{\mathfrak{a}}
\newcommand{\bb}{\mathfrak{b}}
\newcommand{\Ns}{\mathrm{N}}
\newcommand{\Frob}{\operatorname{Frob}}
\newcommand{\Gal}{\operatorname{Gal}}
\newcommand{\Rep}{\operatorname{Rep}}
\newcommand{\Irr}{\operatorname{Irr}}
\newcommand{\Tr}{\operatorname{Tr}}
\newcommand{\HS}{\mathrm{HS}}
\newcommand{\Prob}{\operatorname{Prob}}
\newcommand{\id}{\mathrm{id}}

\PassOptionsToPackage{hyphens}{url}
\usepackage[hidelinks,bookmarks=false]{hyperref}

\title[The Galois-twisted groupoid]{The Galois-twisted arithmetic groupoid:\\ its KMS classification, and why the Tannakian obstruction\\ does not compute it}
\author{Ruqing Chen}
\address{GUT Geoservice Inc., Montreal, Canada}
\email{ruqing@hotmail.com}
\thanks{Sources, code and data for the entire series: \url{https://github.com/Ruqing1963/localized-bost-connes}; this paper is in \texttt{papers/VII-galois-groupoid/}.}
\date{}

\begin{document}

\begin{abstract}
We construct the Galois-twisted arithmetic groupoid $\cG_{G,S}$ over a compact group $G$ and
classify the $\mathrm{KMS}_\beta$ simplex of its $C^*$-algebra completely. The answer is
\[
\mathrm{KMS}_\beta\bigl(\cA_{G,S}\bigr)\ \cong_{\mathrm{affine}}\ \Prob\bigl(E_\beta\backslash G\bigr),
\]
where $E_\beta\le G$ is the essential-value subgroup of the Frobenius cocycle, computed by
the abelian Kakutani criterion inside the closed abelian subgroup
$A=\overline{\Frob(I_S)}$. The right translation action of $G$ is transitive on the extreme
points with stabilizer $E_\beta$.

This is \emph{not} $\Prob(G/N_\beta)$, with $N_\beta$ the Tannakian subgroup of Paper VI
of the series, and the discrepancy is not a technicality. The category
$\cT_\beta=\{\rho:\Xi_\beta(\rho,S)<\infty\}$ determines only the \emph{normal closure} of
$E_\beta$, because $\Rep(G/N)$ sees normal subgroups only, whereas the ergodic decomposition
of a skew product by a non-normal subgroup is a coset space that is not a group. For
$G=S_3$ with $E_\beta$ a transposition subgroup, $\Prob(E_\beta\backslash G)$ is a $2$-simplex
while $\Prob(G/N_\beta)$ is a point: three extremal $\mathrm{KMS}_\beta$ states where the
Tannakian formulation predicts uniqueness. \textbf{The obstruction group $\Xi_\beta$ does not
determine the KMS simplex once $G$ is non-abelian.}

Three natural guesses are corrected. The matrix zero-one law
``$F_\rho F_\rho^*=c_\rho I$'' is false: the cocycle equation conjugates $F_\rho F_\rho^*$, so
ergodicity forces its \emph{spectrum} to be almost everywhere constant and nothing more; for
$G=S_3$, $\rho$ the two-dimensional irreducible and $\lambda=\frac12(\delta_e+\delta_{(12)})$
one gets eigenvalues $\{0,1\}$. The matrix Fourier transform $F_\rho$ does not vanish when
the Frobenius series $\Xi_\beta(\rho,S)$ diverges; it vanishes exactly when $\rho$ has no
$E_\beta$-fixed vector, and this is the refinement of the obstruction data that sees
$E_\beta$. And the ``matrix corner'' route, twisting the coefficient algebra by
$\mathrm{Ad}\,\rho(\Frob_\pp)$, adds nothing: KMS states are tracial on the coefficient
algebra, so inner twists are invisible to them.

Finally we record why the model cannot be genuinely non-abelian: $I_S$ is abelian and
$a\mapsto\Frob_a$ is a homomorphism, so its image is always an abelian subgroup of $G$, the
whole cocycle analysis takes place inside $A$, and $G$ contributes only the coset space. A
genuinely non-abelian theory needs a non-abelian acting semigroup, which by Paper V of the
series costs the high-temperature region where the phase structure lives.
\end{abstract}

\maketitle

\tableofcontents

\section{The Galois-twisted groupoid}

Let $K$ be a number field, $S$ a set of finite primes, $X=\overline{\N}^{\,S}$ with
$\overline\N=\N\cup\{\infty\}$, compact metrizable and totally disconnected, and
$V=\N^S\subset X$ the set of valuation vectors, carrying the forced product measure
$\pi_\beta=\bigotimes_{\pp\in S}\mu_{t_\pp}$, $t_\pp=\Ns\pp^{-\beta}$, $\mu_t(k)=(1-t)t^k$
\cite[Lem.~2.10]{Main}; $\pi_\beta$ charges no hyperplane $\{v_\pp=\infty\}$. Let $\cR_S$ be the
tail equivalence relation on $V$, which is ergodic for $\pi_\beta$ by the Kolmogorov
zero-one law \cite[Lem.~3.1]{Main}. The semigroup $J_S=\N^{(S)}$ acts on $X$ by
$v\mapsto v+a$, with $\infty+a_\pp=\infty$, and $I_S=\Z^{(S)}$ on the dilation
$\widetilde X=\varinjlim(X,+a)$, a locally compact totally disconnected space.

\emph{The series.} The foundational paper is \cite{Main}; Paper I \cite{SeriesI} treats the
factor type and the spectrum of transitions. Papers III, V and VI --- the measure-class
description of $\Xi_\beta$, the free variant, and the Tannakian obstruction for non-abelian
$G$ --- are in preparation; nothing here depends on them, and we refer to them by numeral,
without citation.

\begin{definition}\label{def:groupoid}
Let $G$ be a compact second countable group with Haar measure $m_G$, and let
$\pp\mapsto\Frob_\pp\in G$ be given for $\pp\in S$. Put $\Frob_a=\prod_\pp\Frob_\pp^{a_\pp}$
for $a\in I_S=\bigoplus_{\pp\in S}\Z$, let
\[
\Omega_S=X\times G,\qquad a\cdot(v,g)=\bigl(v+a,\ \Frob_a^{-1}g\bigr),
\]
let $\widetilde\Omega_S=\widetilde X\times G$ be the dilation, and set
\[
\cG_{G,S}=I_S\ltimes\widetilde\Omega_S,\qquad
\cA_{G,S}=\mathbf 1_{\Omega_S}\,C^*\bigl(\cG_{G,S}\bigr)\,\mathbf 1_{\Omega_S}=C(\Omega_S)\rtimes J_S ,
\]
the unital full corner of the groupoid algebra $C_0(\widetilde\Omega_S)\rtimes I_S$, exactly
as $\cA_{K,S}=C(Y_S)\rtimes J_S$ sits in the dilated groupoid algebra in \cite[\S2]{Main}; the
dynamics is $\sigma_t(\mu_\aaa)=\Ns\aaa^{it}\mu_\aaa$, $\sigma_t|_{C(\Omega_S)}=\id$. All
measures below are concentrated on $V\times G$, and we write $(v,g)$ for its points.
\end{definition}

\begin{lemma}[Basic structure]\label{lem:basic}
\begin{enumerate}
\item $a\mapsto\Frob_a$ is a homomorphism $I_S\to G$, so
$A:=\overline{\Frob(I_S)}$ is a closed \emph{abelian} subgroup of $G$, whatever $G$ is.
\item Right translation $R_h(v,g)=(v,gh)$ commutes with the $I_S$-action, so $G$ acts on
$\cA_{G,S}$ by $\sigma$-commuting automorphisms and permutes the $\mathrm{KMS}_\beta$ states.
\item $\cG_{G,S}$ is second countable, étale and amenable, so $C^*_r=C^*$ \cite{Renault}.
\end{enumerate}
\end{lemma}

\begin{proof}
(1) $I_S$ is abelian and $\Frob$ is a homomorphism from it, so the image is an abelian
subgroup, and the closure of an abelian subgroup is abelian. (2) Left and right translations
of $G$ commute. (3) $I_S$ is abelian, hence amenable, acting on a locally compact space by homeomorphisms,
and the corner of a full crossed product by an amenable group is the corner of the reduced
one.
\end{proof}

\begin{remark}[The model cannot be genuinely non-abelian]\label{rem:abelian}
Lemma~\ref{lem:basic}(1) is the fundamental limitation of this construction, and it is worth
stating before any of the analysis. The Frobenius data enter only through the homomorphism
$\Frob:I_S\to G$, whose image is abelian because $I_S$ is. Every cocycle computation below
therefore takes place inside $A$, and the non-abelian group $G$ contributes nothing but the
homogeneous space $A\backslash G$ on which it acts. This is not an artefact of our
normalization: any construction in which an abelian semigroup of ideals acts through
Frobenius elements has the same feature. Genuine non-commutativity would require a
non-abelian acting semigroup; Paper V of the series shows what that costs: partitioning
range projections empty the high-temperature region.
\end{remark}

\section{Disintegration and the matrix Fourier transform}

\begin{proposition}[KMS states are scaling measures]\label{prop:kmsmeasure}
The map $\varphi\mapsto\nu=\varphi|_{C_0(\widetilde\Omega_S)}$ is an affine bijection from
$\mathrm{KMS}_\beta(\cA_{G,S})$ onto the set of Borel probability measures $\nu$ on
$\Omega_S$ with
\begin{equation}\label{eq:scaling}
\nu\bigl(a\cdot E\bigr)=\Ns\aaa^{-\beta}\,\nu(E)\qquad(a\in J_S),
\end{equation}
and every such $\nu$ has $V$-marginal $\pi_\beta$. Disintegrating, $\nu=\pi_\beta\otimes\rho_v$
with $(\rho_v)$ a measurable family in $\Prob(G)$, and \eqref{eq:scaling} is equivalent to
\begin{equation}\label{eq:equivariance}
\rho_{v+e_\pp}=\bigl(\Frob_\pp^{-1}\bigr)_*\rho_v\qquad\pi_\beta\text{-a.e., for every }\pp\in S.
\end{equation}
\end{proposition}

\begin{proof}
As in \cite[Prop.~2.8]{Main}: the canonical expectation onto $C(\Omega_S)$ is
faithful, the KMS condition on $\mu_\aaa f\mu_\aaa^*$ gives \eqref{eq:scaling}, and
conversely each scaling measure reconstructs a KMS state. The $V$-marginal is forced by
\cite[Lem.~2.10]{Main}, whose argument uses only the scaling of the valuation coordinates,
the $G$-direction being measure preserving for the translation.
Splitting \eqref{eq:scaling} into its $V$- and $G$-components gives \eqref{eq:equivariance}.
\end{proof}

\begin{definition}[Matrix Fourier transform]\label{def:F}
For $\rho\in\Irr(G)$ of dimension $d_\rho$ set
\[
F_\rho(v)=\int_G\rho(g)\,d\rho_v(g)\ \in\ M_{d_\rho}(\C),
\]
a contraction, $\|F_\rho(v)\|\le1$, since $\rho_v$ is a probability measure and $\rho$ is
unitary.
\end{definition}

\begin{lemma}[Matrix cohomological equation]\label{lem:coho}
Equation \eqref{eq:equivariance} holds if and only if for every $\rho\in\Irr(G)$
\begin{equation}\label{eq:coho}
F_\rho(v+e_\pp)=\rho(\Frob_\pp)^{-1}\,F_\rho(v)\qquad\pi_\beta\text{-a.e., for every }\pp\in S .
\end{equation}
\end{lemma}

\begin{proof}
If \eqref{eq:equivariance} holds then
\[
F_\rho(v+e_\pp)=\int_G\rho(g)\,d\bigl((\Frob_\pp^{-1})_*\rho_v\bigr)(g)
=\int_G\rho\bigl(\Frob_\pp^{-1}h\bigr)\,d\rho_v(h)
=\rho(\Frob_\pp)^{-1}F_\rho(v).
\]
Conversely, the matrix coefficients of the $\rho\in\Irr(G)$ span a dense subalgebra of $C(G)$
by Peter--Weyl, so \eqref{eq:coho} for all $\rho$ determines \eqref{eq:equivariance}.
\end{proof}

\begin{remark}[The inverse is not cosmetic]\label{rem:inverse}
With the action of Definition~\ref{def:groupoid} the transform satisfies \eqref{eq:coho} with
$\rho(\Frob_\pp)^{-1}$, not $\rho(\Frob_\pp)$. Writing the equation with
$\rho(\Frob_\pp)$ corresponds to the opposite convention $a\cdot(v,g)=(v+a,\Frob_ag)$ and
changes $E_\beta$ into its image under inversion, which for the subgroup $E_\beta$ is
harmless; but in intermediate steps, in particular in the martingale of
Proposition~\ref{prop:gauge}, the two conventions do not agree and mixing them produces a
wrong gauge.
\end{remark}

\section{What the zero-one law says, and what it does not}

\begin{proposition}[Spectral zero-one law]\label{prop:zeroone}
Put $P_\rho(v)=F_\rho(v)F_\rho(v)^*$, positive semidefinite with $\|P_\rho\|\le1$. Then
\eqref{eq:coho} gives
\begin{equation}\label{eq:conj}
P_\rho(v+e_\pp)=\rho(\Frob_\pp)^{-1}P_\rho(v)\,\rho(\Frob_\pp),
\end{equation}
so the \emph{eigenvalues} of $P_\rho(v)$, with multiplicity, are $\cR_S$-invariant and hence
$\pi_\beta$-almost everywhere constant. In particular $\Tr P_\rho(v)$ and $\det P_\rho(v)$
are a.e.\ constant.
\end{proposition}

\begin{proof}
\eqref{eq:conj} is immediate from \eqref{eq:coho} and unitarity. Conjugation preserves the
spectrum, so each elementary symmetric function of the eigenvalues is invariant under the
generators of $\cR_S$, hence a.e.\ constant by ergodicity \cite[Lem.~3.1]{Main}.
\end{proof}

\begin{proposition}[$P_\rho$ need not be scalar]\label{prop:notscalar}
The statement $P_\rho(v)=c_\rho I_{d_\rho}$ a.e.\ is strictly stronger than
Proposition~\ref{prop:zeroone} and is false. Take $G=S_3$, $\rho$ the two-dimensional
irreducible, $\Frob_\pp=e$ for every $\pp\in S$, and
$\rho_v\equiv\lambda:=\tfrac12(\delta_e+\delta_{(12)})$, which satisfies
\eqref{eq:equivariance} trivially. Then
\[
F_\rho=\tfrac12\bigl(I+\rho((12))\bigr),\qquad
P_\rho=\tfrac12\bigl(I+\rho((12))\bigr),
\]
a rank-one projection with eigenvalues $\{0,1\}$, not a multiple of $I_2$.
\end{proposition}

\begin{proof}
$\rho((12))$ is a reflection, so $\rho((12))^*=\rho((12))$ and $\rho((12))^2=I$; hence
$F_\rho^*=F_\rho$ and $F_\rho^2=\tfrac14(I+2\rho((12))+I)=F_\rho$, so $P_\rho=F_\rho F_\rho^*=F_\rho^2=F_\rho$
is a projection, of trace $\tfrac12\Tr(I+\rho((12)))=\tfrac12(2+0)=1$, hence rank one.
\end{proof}

\begin{remark}[Numerical check]\label{rem:numscalar}
In the standard two-dimensional realization $\rho((12))=\begin{psmallmatrix}0&1\\1&0\end{psmallmatrix}$,
one computes $F_\rho=\begin{psmallmatrix}1/2&1/2\\1/2&1/2\end{psmallmatrix}$ and eigenvalues
of $P_\rho$ equal to $\{0,1\}$; conjugating by $\rho((12))$ returns the same spectrum, as
\eqref{eq:conj} requires.
\end{remark}

\begin{proposition}[What the transform sees]\label{prop:gauge}
Let $\rho\in\Irr(G)$ and $U_\pp=\rho(\Frob_\pp)$. The $U_\pp$ commute, lying in $\rho(A)$
with $A=\overline{\Frob(I_S)}$ abelian, so there is a unitary $W$ with
$WU_\pp W^*=\mathrm{diag}(\chi_1(\Frob_\pp),\dots,\chi_{d_\rho}(\Frob_\pp))$ for all $\pp$,
where $\chi_1,\dots,\chi_{d_\rho}\in\widehat A$ are the characters of $A$ occurring in
$\rho|_A$. Let $\Xi_\beta(A)\le\widehat A$ be as in Definition~\ref{def:E} below.
\begin{enumerate}
\item Every measurable solution $F:V\to M_{d_\rho}(\C)$ of \eqref{eq:coho} has the form
$WF(v)=D(v)\,C$ with $C$ a constant matrix whose $r$-th row vanishes unless
$\chi_r\in\Xi_\beta(A)$, and $D(v)=\mathrm{diag}(g_1(v),\dots,g_{d_\rho}(v))$, $g_r$ a fixed
unimodular solution of $g_r(v+e_\pp)=\chi_r(\Frob_\pp)^{-1}g_r(v)$ when $\chi_r\in\Xi_\beta(A)$
and $g_r\equiv1$ otherwise.
\item \eqref{eq:coho} has a unitary solution $W_\rho:V\to U(d_\rho)$, unique up to right
multiplication by a constant unitary, if and only if every $\chi_r$ lies in $\Xi_\beta(A)$,
i.e.\ $\Xi_\beta(\rho,S)<\infty$, i.e.\ $\rho|_{E_\beta}=1$.
\item \eqref{eq:coho} has a nonzero solution if and only if some $\chi_r$ lies in
$\Xi_\beta(A)$, i.e.\ $\rho$ has a nonzero $E_\beta$-fixed vector. In particular
$\Xi_\beta(\rho,S)=\infty$ does \emph{not} force $F_\rho\equiv0$: for $G=S_3$,
$\Frob_\pp=(12)$ for all $\pp$, $\zeta_{K,S}(\beta)=\infty$ and $\rho$ the two-dimensional
irreducible, $\Xi_\beta(\rho,S)=4\log\zeta_{K,S}(\beta)+O(1)=\infty$, while the constant
family $\rho_v\equiv\frac12(\delta_e+\delta_{(12)})$ is equivariant with
$F_\rho\equiv\frac12(I+\rho((12)))\ne0$.
\end{enumerate}
\end{proposition}

\begin{proof}
Put $H=WF$; then \eqref{eq:coho} reads $H_{rs}(v+e_\pp)=\chi_r(\Frob_\pp)^{-1}H_{rs}(v)$ for
each entry. If $\chi_r\notin\Xi_\beta(A)$, every solution of this scalar equation vanishes
a.e.: by \cite[Lem.~3.1]{Main} $|H_{rs}|$ is a.e.\ constant, and if the constant were
positive, \cite[Prop.~3.3]{Main}, whose proof uses only $|\chi_r(\Frob_\pp)|=1$, would make
$\sum_\pp\Ns\pp^{-\beta}|1-\chi_r(\Frob_\pp)|^2$ converge. If $\chi_r\in\Xi_\beta(A)$,
\cite[Prop.~3.4]{Main} supplies a unimodular solution $g_r$, and $H_{rs}/g_r$ is
$\cR_S$-invariant, hence a.e.\ constant. This is (1). For (2), a unitary solution needs all
rows of $C$ nonzero, hence all $\chi_r\in\Xi_\beta(A)$; conversely $W^*D(v)$ is then a
unitary solution, and two unitary solutions differ by an invariant, hence constant, unitary.
The equivalence with $\Xi_\beta(\rho,S)<\infty$ is $\|I-U_\pp\|_{\HS}^2=\sum_r|1-\chi_r(\Frob_\pp)|^2$,
and with $\rho|_{E_\beta}=1$ is Pontryagin duality, $\Xi_\beta(A)=E_\beta^\perp$ since
$\widehat A$ is discrete. (3) is read off from (1) in the same way; the $E_\beta$-fixed
vectors of $\rho$ are the sum of the eigenspaces of those $\chi_r$ that are trivial on
$E_\beta$. In the example, $\rho((12))$ is a reflection, so the two characters of
$A=\{e,(12)\}$ occurring in $\rho|_A$ are the trivial one and the sign, and the
displayed family satisfies $\rho_{v+e_\pp}=((12))_*\rho_v=\rho_v$.
\end{proof}

\begin{remark}[The refinement that sees $E_\beta$]\label{rem:fixed}
The Dirichlet series $\Xi_\beta(\rho,S)$ records whether $\rho$ is trivial on $E_\beta$.
The vanishing or not of the matrix Fourier transform records whether $\rho^{E_\beta}\ne0$;
and the a.e.\ constant rank of $P_\rho=F_\rho F_\rho^*$ from Proposition~\ref{prop:zeroone},
for the extremal equivariant families of Theorem~\ref{thm:main} below, is
$\dim\rho^{E_\beta}$: for such a family $\rho_v$ is Haar measure on a coset $y(v)E_\beta h$,
and $F_\rho(v)=\rho(y(v))\,P^{E_\beta}\rho(h)$ with $P^{E_\beta}$ the projection onto
$\rho^{E_\beta}$. The function $\rho\mapsto\dim\rho^{E_\beta}$ on $\Irr(G)$ is the multiplicity function
of the permutation representation $\mathrm{Ind}_{E_\beta}^G\mathbf 1$ on $L^2(E_\beta\backslash G)$
\cite{Mackey}: it determines $|E_\beta\backslash G|=\sum_\rho d_\rho\dim\rho^{E_\beta}$, the
number of extremal $\mathrm{KMS}_\beta$ states, and the decomposition of that permutation
representation, which distinguishes $E_\beta$ from its normal closure --- in
Example~\ref{ex:S3} it is $\mathbf 1\oplus\rho_2$ rather than $\mathbf 1$. It does not in
general determine $E_\beta$ up to conjugacy, since non-conjugate subgroups can have the
same permutation character; but it sees the simplex, which the Frobenius series does not.
\end{remark}

\section{The KMS classification}

\begin{definition}\label{def:E}
Let $A=\overline{\Frob(I_S)}$, abelian by Lemma~\ref{lem:basic}(1), and let
\[
\Xi_\beta(A):=\Bigl\{\chi\in\widehat A:\ \sum_{\pp\in S}\Ns\pp^{-\beta}\bigl|1-\chi(\Frob_\pp)\bigr|^2<\infty\Bigr\},
\qquad
E_\beta:=\Xi_\beta(A)^\perp\le A ,
\]
the essential-value subgroup of the Frobenius cocycle $c(v+a,v)=\Frob_a$ in the sense of
Schmidt \cite{Schmidt}; for the compact target $A$ the essential values of $c$ are exactly
the annihilator of the coboundary characters, which is how $E_\beta$ enters below.
\end{definition}

\begin{theorem}[Classification]\label{thm:main}
The map $\varphi\mapsto\nu$ of Proposition~\ref{prop:kmsmeasure} induces an affine
isomorphism
\begin{equation}\label{eq:main}
\mathrm{KMS}_\beta\bigl(\cA_{G,S}\bigr)\ \cong\ \Prob\bigl(E_\beta\backslash G\bigr),
\end{equation}
and the right translation action of $G$ is transitive on the extreme points, with the
stabilizer of the base point equal to $E_\beta$. In particular the $\mathrm{KMS}_\beta$
state is unique if and only if $E_\beta=G$, which for $G$ non-abelian forces $G$ itself to be
abelian; so for non-abelian $G$ the simplex is never a point.
\end{theorem}

\begin{proof}
By Proposition~\ref{prop:kmsmeasure} we must describe the equivariant families $(\rho_v)$.
The $I_S$-action on $\Omega_S$ is the skew product of $\cR_S$ with the constant $G$-valued
cocycle $c(v+a,v)=\Frob_a$, whose values lie in $A$. Extreme equivariant families correspond
to ergodic components of that skew product \cite{FM}.

Decompose $G$ into the cosets $Ag$; the $I_S$-action preserves each, acting on
$Ag\cong A$ by translation by $\Frob_a^{-1}$. An equivariant family disintegrates over
$A\backslash G$ into equivariant families on the cosets, so the extremal ones are
supported on a single coset $Ag$ and are extremal there. Within one coset the problem is
that of \cite[Thm.~3.6]{Main} with $A$ in place of $G_S$ and $\Frob_\pp^{-1}$ in place of
$\sigma_\pp$: the Fourier coefficients $\hat\rho_v(\chi)$, $\chi\in\widehat A$, satisfy the
scalar cohomological equation with the unimodular constants $\chi(\Frob_\pp)^{\mp1}$, so
by \cite[Lem.~3.1, Props.~3.3, 3.4]{Main} they vanish unless $\chi\in\Xi_\beta(A)$, and the
extremal families are the Haar measures on the cosets of $E_\beta=\Xi_\beta(A)^\perp$ in
$A$ moving by $\Frob_\pp^{-1}$. Transported to $Ag$, these are the Haar measures on the
subsets $aE_\beta g=E_\beta\,ag$, $a\in A$. Hence the extremal equivariant families are
indexed by the right cosets $E_\beta h$, $h\in G$, and the simplex of equivariant families
is $\Prob(E_\beta\backslash G)$, giving \eqref{eq:main}.

Right translation by $h\in G$ carries the component of $E_\beta g$ to that of $E_\beta gh$,
so the action is transitive, with stabilizer of $E_\beta e$ equal to $E_\beta$. Uniqueness
means $|E_\beta\backslash G|=1$, i.e.\ $E_\beta=G$; but $E_\beta\le A$ is abelian.
\end{proof}

\begin{corollary}[Abelian case]\label{cor:abelian}
If $G$ is abelian and $\Frob(I_S)$ is dense, then $A=G$, $E_\beta=\Xi_\beta^\perp$ and
\eqref{eq:main} is $\Prob(G/\Xi_\beta^\perp)$, recovering \cite[Thm.~3.6]{Main}.
\end{corollary}

\section{The Tannakian obstruction does not compute the simplex}

\begin{theorem}[Discrepancy]\label{thm:discrepancy}
Let $\Xi_\beta(\rho,S)=\sum_{\pp\in S}\Ns\pp^{-\beta}\|I-\rho(\Frob_\pp)\|_{\HS}^2$, let
$\cT_\beta(S)=\{\rho:\Xi_\beta(\rho,S)<\infty\}$, and let $N_\beta=\bigcap_{\rho\in\cT_\beta}\ker\rho$
be the Tannakian subgroup of Paper VI, for which $\cT_\beta(S)=\Rep(G/N_\beta)$. Then
\[
\cT_\beta(S)=\bigl\{\rho:\rho|_{E_\beta}=1\bigr\},
\qquad
N_\beta=\text{the normal closure of }E_\beta\text{ in }G .
\]
Consequently $\Prob(E_\beta\backslash G)$ and $\Prob(G/N_\beta)$ agree if and only if
$E_\beta$ is normal in $G$, and in general
\[
\bigl|E_\beta\backslash G\bigr|\ \ge\ \bigl|G/N_\beta\bigr| ,
\]
with strict inequality whenever $E_\beta$ is not normal. \textbf{The obstruction group
$\Xi_\beta$ therefore does not determine the $\mathrm{KMS}_\beta$ simplex when $G$ is
non-abelian.}
\end{theorem}

\begin{proof}
By Proposition~\ref{prop:gauge}(2), $\Xi_\beta(\rho,S)<\infty$ iff $\rho|_{E_\beta}=1$, for
irreducible $\rho$ and hence for all $\rho$ by additivity of $\Xi_\beta$ over direct sums.
A representation is
trivial on $E_\beta$ iff it is trivial on the normal closure $\langle\langle E_\beta\rangle\rangle$,
since $\ker\rho$ is normal; hence $\cT_\beta=\Rep(G/\langle\langle E_\beta\rangle\rangle)$ and
$N_\beta=\langle\langle E_\beta\rangle\rangle$. Since $E_\beta\subseteq N_\beta$, the natural
surjection $E_\beta\backslash G\to N_\beta\backslash G=G/N_\beta$ gives the inequality, and it
is a bijection iff $E_\beta=N_\beta$, i.e.\ iff $E_\beta$ is normal.
\end{proof}

\begin{example}[Three states where the Tannakian prediction is one]\label{ex:S3}
Let $G=S_3$, $\Frob_\pp=(12)$ for every $\pp\in S$, and $\beta$ with
$\zeta_{K,S}(\beta)=\infty$. Then $A=E_\beta=\{e,(12)\}$: indeed the nontrivial character
$\chi$ of $A$ has $\chi(\Frob_\pp)=-1$ for all $\pp$, so
$\sum\Ns\pp^{-\beta}|1-\chi(\Frob_\pp)|^2=4\zeta_{K,S}(\beta)=\infty$ and
$\Xi_\beta(A)=\{1\}$.
\begin{itemize}
\item By Theorem~\ref{thm:main} the simplex is $\Prob(E_\beta\backslash S_3)$, a
$2$-simplex: \textbf{three} extremal $\mathrm{KMS}_\beta$ states, permuted transitively by
$S_3$ with stabilizer of order $2$.
\item By Theorem~\ref{thm:discrepancy}, $\cT_\beta$ consists of the representations trivial
on $(12)$. The sign character has $\mathrm{sgn}((12))=-1$ and the two-dimensional
irreducible restricts to $E_\beta$ with eigenvalues $\pm1$; both are nontrivial. So
$\cT_\beta=\{\mathbf 1\}$, $N_\beta=S_3$, and $\Prob(G/N_\beta)$ is a \textbf{point}.
\end{itemize}
Three against one. The Tannakian formulation predicts uniqueness where the system has
spontaneous symmetry breaking with a three-element phase space.
\end{example}

\begin{remark}[Why the two disagree]\label{rem:why}
Tannaka--Krein reconstructs a group from a representation category, and a representation
category can only see \emph{normal} subgroups, being determined by kernels. The ergodic
decomposition of a skew product by a subgroup $E$ is the coset space $E\backslash G$, which
for non-normal $E$ is a homogeneous space carrying no group structure. In the abelian case
every subgroup is normal and the distinction is invisible, which is why the $GL_1$ theory
gave no warning. The correct reading of the Tannakian theorem of Paper VI is therefore
that $\cT_\beta$ computes the largest \emph{quotient group} through which the obstruction
factors;
that quotient is a coarsening of the true invariant, which is the pair $(G,E_\beta)$ up to
conjugacy, equivalently the transitive $G$-space $E_\beta\backslash G$.
\end{remark}

\begin{corollary}[What is and is not an invariant]\label{cor:invariant}
The $\mathrm{KMS}_\beta$ simplex of $\cA_{G,S}$, together with its $G$-action, determines
$E_\beta$ up to conjugacy. The Dirichlet obstruction data
$\{\Xi_\beta(\rho,S)\}_{\rho\in\Irr(G)}$ determine only $N_\beta$. Hence two systems with
non-conjugate $E_\beta$ having the same normal closure have identical obstruction data and
non-isomorphic simplices.
\end{corollary}

\section{Route B: the matrix corner}

\begin{proposition}[Inner twists are invisible]\label{prop:corner}
Let $\rho:G\to U(n)$ be any unitary representation, let $(B,\alpha)$ be a unital
$C^*$-algebra with an action of $J_S$ by injective endomorphisms of the kind in
\cite[\S2]{Main}, and let $\alpha_\rho$ be the twisted action on $B\otimes M_n(\C)$,
\[
\alpha_{\rho,\pp}(b\otimes m)=\alpha_\pp(b)\otimes\rho(\Frob_\pp)^{-1}m\,\rho(\Frob_\pp).
\]
Give $B_\rho=(B\otimes M_n(\C))\rtimes_{\alpha_\rho}J_S$ the dynamics fixing $B\otimes M_n$
and scaling $\mu_\aaa$ by $\Ns\aaa^{it}$. Then every $\mathrm{KMS}_\beta$ state of $B_\rho$
restricts to $B\otimes M_n$ as $\varphi_0\otimes\tfrac1n\Tr$ with $\varphi_0$ the
restriction of a $\mathrm{KMS}_\beta$ state of $B\rtimes_\alpha J_S$, and
$\varphi\mapsto\varphi_0$ is an affine bijection
$\mathrm{KMS}_\beta(B_\rho)\cong\mathrm{KMS}_\beta(B\rtimes_\alpha J_S)$. The twist by
$\rho$ does not change the simplex.
\end{proposition}

\begin{proof}
For $x,y$ in the fixed-point algebra of the dynamics, which contains $B\otimes M_n$, the KMS
condition gives $\varphi(xy)=\varphi(y\sigma_{i\beta}(x))=\varphi(yx)$; so $\varphi$ is
tracial on $B\otimes M_n$, hence of the form $\varphi_0\otimes\frac1n\Tr$. A KMS state
vanishes on $\mu_\aaa x\mu_\bb^*$ for $\aaa\ne\bb$, and on $\mu_\aaa x\mu_\aaa^*$ it
equals $\Ns\aaa^{-\beta}\varphi(x)$; so $\varphi$ is determined by its restriction, and
the KMS condition on $B_\rho$ reduces to $\varphi\circ\alpha_{\rho,\pp}=\Ns\pp^{-\beta}\varphi$
on $B\otimes M_n$. Since $\Tr\bigl(\rho(\Frob_\pp)^{-1}m\rho(\Frob_\pp)\bigr)=\Tr m$, this
reads $\varphi_0\circ\alpha_\pp=\Ns\pp^{-\beta}\varphi_0$, the KMS condition of the untwisted
system, and conversely every such $\varphi_0$ gives a KMS state $\varphi_0\otimes\frac1n\Tr$
of $B_\rho$ by the same computation.
\end{proof}

\begin{remark}\label{rem:routeB}
``Route B'' proposed to import the non-abelian representation through a matrix-valued
coefficient algebra $C(Y_S)\otimes M_n(\C)$ twisted by conjugation with $\rho(\Frob_\pp)$.
Proposition~\ref{prop:corner} shows that this adds nothing, for a reason more basic than
the spectral zero-one law: equilibrium states are tracial on the coefficient algebra, and
a trace does not see inner automorphisms. The representation can only enter through the
base space, which is Route A. (The twisted algebra is related to $\cA_{G,S}$: the crossed
product of $\cA_{G,S}$ by the right translation action of $G$ decomposes, by Peter--Weyl
and the Stone--von Neumann theorem, into blocks indexed by $\Irr(G)$, and the $\rho$-block
is Morita equivalent to $B_\rho$ with $B=C(X)$. That description is consistent with the
proposition, and we do not use it.)
\end{remark}

\section{Assessment}

\[
\begin{array}{lll}
\text{Statement} & \text{Status} & \text{Reference}\\\hline
\text{groupoid }\cG_{G,S}\text{ well defined, }C^*_r=C^*&\text{true}&\text{Lem.~1.2}\\
\text{KMS states}\leftrightarrow\text{scaling measures}&\text{true}&\text{Prop.~2.1}\\
F_\rho(v+e_\pp)=\rho(\Frob_\pp)F_\rho(v)&\text{inverse missing}&\text{Lem.~2.3}\\
F_\rho F_\rho^*=c_\rho I&\text{false}&\text{Prop.~3.2}\\
\text{spectrum of }F_\rho F_\rho^*\text{ a.e.\ constant}&\text{true}&\text{Prop.~3.1}\\
\Xi_\beta(\rho,S)=\infty\Rightarrow F_\rho=0&\text{false}&\text{Prop.~\ref{prop:gauge}(3)}\\
F_\rho\ne0\iff\rho^{E_\beta}\ne0&\text{true}&\text{Prop.~\ref{prop:gauge}, Rem.~\ref{rem:fixed}}\\
\text{unitary gauge exists}\iff\Xi_\beta<\infty&\text{true}&\text{Prop.~\ref{prop:gauge}(2)}\\
\mathrm{KMS}_\beta\cong\Prob(G/N_\beta)&\textbf{false}&\text{Ex.~5.2}\\
\mathrm{KMS}_\beta\cong\Prob(E_\beta\backslash G)&\textbf{true}&\text{Thm.~4.2}\\
\Xi_\beta\text{ determines the simplex}&\text{false for non-abelian }G&\text{Cor.~5.4}\\
\text{Route B adds information}&\text{no: inner twists are invisible}&\text{Prop.~6.1}\\
\text{the model is genuinely non-abelian}&\text{no}&\text{Rem.~1.3}
\end{array}
\]

The construction succeeds and the classification is complete, so the question left open in
Paper VI is answered. The answer is not the expected one, in two ways. The simplex
is a space of probability measures on a homogeneous space rather than on a group, and the
Dirichlet obstruction data --- which is all that the $L$-functions of Paper VI can
see --- determine that homogeneous space only up to replacing $E_\beta$ by its normal
closure. Any programme that hopes to read a non-abelian phase structure off automorphic
$L$-functions must therefore contend with Corollary~\ref{cor:invariant}: the $L$-functions
know $N_\beta$, and $N_\beta$ is not enough.

Three questions. First, Remark~\ref{rem:fixed} identifies a refinement of the
obstruction data, $\rho\mapsto\dim\rho^{E_\beta}$, read off from the rank of
$F_\rho F_\rho^*$, which sees the size of the simplex and the permutation character of
$E_\beta\backslash G$; whether this function of $\rho$ has an expression in $L$-functions,
as $\Xi_\beta(\rho,S)$ has in Paper VI, and what further datum pins down $E_\beta$ up to
conjugacy, we do not know.
Second, by
Remark~\ref{rem:abelian} no construction with an abelian acting semigroup can produce a
non-abelian $E_\beta$; the interesting case $E_\beta$ non-normal already arises with
$E_\beta$ abelian and $G$ non-abelian, as in Example~\ref{ex:S3}, but a theory in which the
essential-value group itself is non-abelian would need the acting semigroup to be
non-abelian, and Paper V shows that this empties the high-temperature region. Whether some intermediate structure --- an Ore semigroup that is non-abelian but has
nested rather than partitioning range projections --- exists is open. Third, we have not
determined which pairs $(G,E_\beta)$ are realizable by an actual arithmetic
$\pp\mapsto\Frob_\pp$; Example~\ref{ex:S3} uses a constant assignment, and Chebotarev should
allow much more, along the lines of \cite[Thm.~5.1]{SeriesI}.

\begin{thebibliography}{9}
\bibitem{FM} J. Feldman, C. C. Moore, \emph{Ergodic equivalence relations, cohomology, and von Neumann algebras II}, Trans. Amer. Math. Soc. 234 (1977), 325--359.
\bibitem{Mackey} G. W. Mackey, \emph{Ergodic theory and virtual groups}, Math. Ann. 166 (1966), 187--207.
\bibitem{Renault} J. Renault, \emph{A Groupoid Approach to $C^*$-Algebras}, Lecture Notes in Math. 793, Springer, 1980.
\bibitem{Schmidt} K. Schmidt, \emph{Cocycles on Ergodic Transformation Groups}, Macmillan Lectures in Math. 1, Delhi, 1977.
\bibitem{Main} R. Chen, \emph{KMS state uniqueness and phase transitions in localized Bost--Connes systems}, preprint, 2026, \newline\url{https://doi.org/10.5281/zenodo.22152101}.
\bibitem{SeriesI} R. Chen, \emph{Localized Bost--Connes systems I: factor type and the spectrum of transitions}, preprint, 2026, DOI \href{https://doi.org/10.5281/zenodo.22160827}{10.5281/zenodo.22160827}.
\end{thebibliography}

\end{document}
